The base field throughout the paper is $\C$.
\subsection{Notation}
\begin{nota}
 Let $G$ be a finite group.
 \begin{itemize}
 \item We denote by $1_G$ its unit element, by $Z(G)$ its center, by $G//G$ the set of conjugacy classes of $G$, and by $C(g)$ the centralizer of $g\in G$.
  \item We denote by $\mathrm{Irr}(G)$ the set of its irreducible representations and by $\triv$ its trivial representation. 
  \item We denote by $\C[G]^{adj}$ the conjugation representation of $G$ on its group algebra.
  \item For any representation $\rho$ of $G$, we denote by $\chi_\rho: G\to \C$ its character, where $\chi_\rho(g)\coloneqq tr (\rho(g))$. By abuse of notation, we use the same notation for the corresponding function $\chi_\rho: G//G \to \C$.
 \item We say that a group $G$ is $p$-torsion, where $p$ is prime, if $g^p = 1_G$ for each $g \in G$.
 \item Given a subgroup $H < G$, a $G$-representation $\rho$ and an $H$-representation $\tau$, we denote by $\rho \downarrow_H$ the restriction of $\rho$ to $H$ and by $\tau \uparrow^G_H$ the induction of $\tau$ to $G$.
 \end{itemize}

\end{nota}

\begin{nota}
 {For $n\geq 2$, we will denote by $C_n$ the cyclic group of order $n$. For prime $p$, the elementary abelian group $p$-group $C_p^{\times k}$ will sometimes be denoted $\mathbb{F}_p^k$.}
\end{nota}

\begin{nota}
 {For $n\geq 3$,} we will denote by $\mathrm{Dih}_n$ the dihedral group of order $2n$ (group of symmetries of the regular $n$-gon). {We will also denote by $\mathrm{Dih}_2$ the group $C_2 \times C_2$ (considered as the group of symmetries of a digon).}
\end{nota}

\begin{nota}
 Let $\Gamma$ be a directed graph.  
 \begin{itemize}
 \item We denote by $X(\Gamma)$ the set of vertices of $\Gamma$. 
 \item We denote by $\mathtt{N}(x)$ the multiset of neighbors of a vertex $x$ in $\Gamma$: those are $y\in X(\Gamma)$ with an edge $x\to y$ in $\Gamma$ (if there are several edges $x \to y$, then $y$ appears with appropriate multiplicity in $\mathtt{N}(x)$.
 \item A loop is an edge connecting a vertex to itself). 
 \item A graph $\Gamma$ is called {\it simply-laced} if for each $x, y \in X(\Gamma)$ it contains at most one edge $x \to y$.
 \item A {\it walk} in $\Gamma$ is a sequence of edges $(e_0, \ldots, e_s)$ such that the end of $e_i$ is the beginning of $e_{i+1}$ for each $i<s$ (repetitions of edges are allowed!). A walk is called a {\it circuit} if the end of $e_s$ is the beginning of $e_{0}$. A {\it path} is a walk without repetitions of edges.
 \item Given a pair of edges in opposite directions (or a loop), we will usually draw a single {\it undirected} edge instead. If all the edges in the graph are undirected, we will say that this graph is {\it undirected}. 
 \item We say that $\Gamma$ is an {\it undirected tree} if $\Gamma$ is undirected and when considered as an undirected graph, it contains no circuits (in particular, it is simply-laced and without loops!).
 \end{itemize}
 
\end{nota}
\begin{nota}
 Given a graph $\Gamma$ with vertex set $X$, we denote by $A_{\Gamma}$ the {\it adjacency matrix} of this graph. That is, $A_{\Gamma}$ is an $\abs{X} \times \abs{X}$ matrix with integer coefficients whose rows and columns are indexed by the set $X$. The $(x,y) \in X \times X$ entry of $A_{\Gamma}$ is the number of edges in $\Gamma$ from $x$ to $y$. The matrix $A_{\Gamma}$ naturally defines a linear endomorphism of the space of functions $X \to \C$, and we will identify $A_{\Gamma}$ with this endomorphism. 
\end{nota}

\begin{rema}
 A graph $\Gamma$ is undirected if and only if $A_{\Gamma}$ is symmetric.
\end{rema}



\subsection{Preliminaries on extra special groups}\label{ssec:extra_special}
A special class of finite groups which appear in this paper are extra special groups. For an introduction to extra special groups, see \cite{Aschbacher} and the appendix of \cite{Diaconis}.

\begin{defi}
 Let $p$ be a prime number. A finite group $G$ is called an {\it extra special $p$-group} if $Z(G) \cong C_p$ and $G/Z(G) \cong (C_p)^N {= \F_p^N}$ for some $N\geq 0$. 
\end{defi}
\begin{exam}\label{ex:extra_special_groups}
\mbox{}
\begin{itemize}
 \item There is just one extra special group of order $p$, up to isomorphism, which is~$C_p$. 
 \item There are two examples of extra special $p$-groups of order $p^3$: the 
Heisenberg group $\mathrm{Heis}(p)$ given by upper-triangular $3\times 3$ matrices over the field 
$\mathbb{F}_p$ with $1$'s on the diagonal, and the semidirect product $C_{p^2} 
\rtimes C_p$ where $C_p$ acts non-trivially on $C_{p^2}$. These are non-isomorphic for $p>2$. 
 \item For $p=2$ the above examples coincide: we have $\mathrm{Heis}(2)\cong C_4 \rtimes C_2 \cong \mathrm{Dih}_4$. Another (non-isomorphic) extra special group of 
order $8$ is the quaternion group $Q_8$.
 \item Given a finite-dimensional vector space $V$ over the field $\F_p$, we can 
define the generalized Heisenberg group $\mathrm{Heis}(V)$ as follows: the underlying set 
is $V \times V^* \times \F_p$ and the multiplication is given by $$(\vec{a}, 
\vec{b}, c)(\vec{a'}, \vec{b'}, c') = (\vec{a} +\vec{a}, \vec{b}+\vec{b'}, c + 
c' + \vec{b'}\cdot\vec{a})$$ where $-\cdot -$ denotes the obvious pairing $V^* 
\times V \to \F_p$. This group $\mathrm{Heis}(V)$ is an extra special $p$-group of order 
$p^{1+2\dim V}$.
\end{itemize}

 
\end{exam}

It turns out that such groups exist only if $N$ is even, i.e. $N=2n$ for some 
$n\geq 0$; moreover, for each power $p^{1+2n}$ ($n\geq 1$) there exist precisely 
two isomorphism classes of extra special $p$-groups of order $p^{1+2n}$. 

For $p>2$, any extra special $p$-group of order at least $p^3$ is a central 
product of several copies of the two groups of order $p^3$ appearing in Example 
\ref{ex:extra_special_groups}. If all of the $n$ factors appearing in the 
central product are isomorphic to $\mathrm{Heis}(p)$ then the obtained group is 
generalized Heisenberg group $\mathrm{Heis}(\F_p^n)$ and it is $p$-torsion. If at least 
one of the $n$ factors is of the form $C_{p^2} \rtimes C_p$ then we obtain 
another, non-isomorphic extra special group of order $p^{1+2n}$, but now having 
some elements of order~$p^2$.

For $p=2$, a similar situation occurs: any extra special $2$-group of order at 
least $8$ is a central product of several copies of $\mathrm{Dih}_4$ and $Q_8$. If 
all of the $n$ factors appearing in the central product are isomorphic to 
$\mathrm{Dih}_4$ then the obtained group is the generalized Heisenberg group 
$\mathrm{Heis}(\F_2^n)$. If at least one of the $n$ factors is of the form $Q_8$ then we 
obtain another, non-isomorphic extra special group of order $2^{1+2n}$.

An extra special $p$-group $G$ of order $p^{1+2n}$ has precisely $p^{2n}$ 
representations of dimension $1$ on which the center $Z(G)$ acts trivially, and 
$p-1$ representations of dimension $p^n$ on which the center $Z(G)$ acts by a 
non-trivial character.

\subsection{Preliminaries on McKay graphs}
The definitions and statements of this subsection are taken from \cite{McKay}.
\begin{defi}[McKay graph]
Let $G$ be a finite group, $\rho$ a complex representation of~$G$. The 
{\it McKay graph} $\Gamma(G, \rho)$ (sometimes called ``McKay quiver'') for the 
pair $(G, \rho)$ has vertex set $\mathrm{Irr}(G)$; the number of edges $\mu \to \lam$, 
where $\mu,\lam \in \mathrm{Irr}(G)$, is 
$$\dim \Hom_G(\mu \otimes \rho,\lam) = [\mu\otimes \rho:\lam]_G$$ where the 
latter denotes the multiplicity of $\lam$ in the $G$-representation $\mu\otimes 
\rho$.
\end{defi}
When drawing a McKay graph, we will usually mark by a $\bigstar$ symbol the vertex corresponding to the trivial representation. The other vertices are each marked by the dimension of the corresponding irreducible representation. 

\begin{exam}

\mbox{}

 \begin{enumerate}
 \item Let $G= C_2$ and $\sgn$ be the non-trivial irreducible representation 
of $C_2$ (``sign representation''). 
 Then $\Gamma(C_2, \sgn)$ is an undirected loop-less graph on $2$ vertices:
 $$
\begin{tikzcd}[arrows=-]
 \bigstar\arrow[r] & \circled{1}
\end{tikzcd}
$$

 \item Let $G= \Z/n\Z$ and $\tau$ be its non-trivial irreducible 
representation. 
Then $\Gamma(\Z/n\Z, \tau)$ is a directed cycle graph on $n$ vertices, without 
double edges nor loops. 
 \item {Let $n\geq 1$ and let $G = \mathbb{F}_2^n = C_2^{\times n}$ be the corresponding elementary abelian $2$-group}. Let $\rho = \bigoplus_{0 \leq s 
\leq n-1} \triv^{\boxtimes s} \boxtimes \sgn \boxtimes \triv^{\boxtimes n-s-1}
  $. The McKay graph $\Gamma(\mathbb{F}_2^n, \rho)$ is then the $n$-dimensional cube 
$\{0, 1\}^n$ with undirected edges $\xymatrix{x \ar@{-}[r]&y}$ whenever $x,y \in \{0, 1\}^n$ differ by 
precisely one coordinate.
 \item Let $G = S_3 = \mathrm{Dih}_3$ be the symmetric group on $3$ letters. We denote its 
irreducible representations by $\triv, \sgn, \mathtt{ref}$ where $\sgn$ is the 
$1$-dimensional sign representation and $\mathtt{ref}$ is the two-dimensional reflection 
representation $\{(x_1, x_2, x_3) \in \C^3, \sum x_i =0\}$. The McKay graph 
$\Gamma(S_3, \mathtt{ref})$ is 
$$ \begin{tikzcd}[arrows=-]
 \bigstar\arrow[r] & \circled{2} \arrow[ld] \arrow[loop right]{r}\\
 \circled{1}
\end{tikzcd}
$$
 This is an undirected graph on $3$ vertices, with a single loop.
 \item Consider the dihedral group $\mathrm{Dih}_{4}$ and let $\tau$ be its tautological $2$-dimensional (irreducible) 
representation. The McKay graph of $(\mathrm{Dih}_{4}, \tau)$ looks as follows:
$$
\begin{tikzcd}[arrows=-]
    \bigstar&& \circled{1} \\
&\ \circled{2} \arrow[ru]\arrow[lu]\arrow[ld]\arrow[rd]\\ 
 \circled{1} && \circled{1} 
\end{tikzcd}
$$
Here the vertex in the center corresponds to $\tau$.
 This graph coincides with the McKay graph for $(Q_8, \mu)$ where $Q_{8}$ is the 
quaternion group and $\mu$ is the obvious $2$-dimensional (irreducible) 
representation given by $\mathbb{H} \cong \C^2$. This coincidence is not 
surprising, given that the character tables of $\mathrm{Dih}_{4}$, $Q_8$ coincide.
 \end{enumerate}
\end{exam}

\begin{exam}\label{ex:extra_special}
Consider an extra-special $2$-group $G$ of order $2^{1+2n}$, $n\geq 0$ as 
described in Section~\ref{ssec:extra_special} (for $n=0$, the corresponding extra special group is unique, and it is just $C_2$). Then $G$ has $2^{2n}+1$ 
irreducible representations: $2^{2n}$ irreducible representations of dimension 
$1$ coming from $G/Z(G) \cong \mathbb{F}_2^{2n}$ and one irreducible representation 
$\rho$ of dimension $2^n$. The McKay graph for $(G, \rho)$ is an unoriented 
connected graph of th following form: it has one vertex $\rho$ with ${4}^n$ 
unoriented edges coming out, and ${4}^n$ vertices connected only to $\rho$ by a 
single unoriented edge. {We will call such a graph {\it a hedgehog with $4^n$ spines}.} Below is an example for $n=2$: 
$$
\begin{tikzcd}[arrows=-]
    \bigstar& \circled{1}& \circled{1}& \circled{1}& \circled{1}\\
   \circled{1} &&&& \circled{1} \\ 
\circled{1} & & \circled{\normalsize $2^n$} \arrow[rr]\arrow[rru]\arrow[rruu]\arrow[ruu]\arrow[uu]\arrow[ll]
\arrow[llu]\arrow[lluu]\arrow[luu]\arrow[rrd]\arrow[rrdd]\arrow[rdd]\arrow[dd]
\arrow[lld]\arrow[lldd]\arrow[ldd]&& \circled{1} \\
     \circled{1} &&&& \circled{1} \\
     \circled{1}& \circled{1}& \circled{1}& \circled{1}& \circled{1}\\
\end{tikzcd}
$$
\end{exam}

The following statement is well-known (see \cite[Proposition 2]{McKay}):
\begin{lemma}\label{lem:undirected}
 Let $\rho$ be a representation of the group $G$. Then $\Gamma(G, \rho^*)$ is 
obtained from $\Gamma(G, \rho)$ by a reversal of arrows (i.e. $A_{\Gamma(G, 
\rho^*)} = A_{\Gamma(G, \rho)}^T$). In particular, $\rho\cong \rho^*$ if and 
only if $\Gamma(G, \rho)$ is undirected (i.e. $A_{\Gamma(G, \rho)}$ is a 
symmetric matrix).
\end{lemma}


The spectrum of the adjacency matrix of a McKay graph has an explicit description, see 
\cite[Proposition 6]{McKay}, \cite[Proposition 2.3]{Browne}, and 
\cite{Steinberg}:

\begin{prop}\label{prop:spectra_adj_matrix}
 The eigenvalues of $A_{\Gamma(G, \rho)}$ form the multiset $$\{\chi_{\rho}(g) |g\in G//G\}.$$ A corresponding eigenvector for the 
eigenvalue $\chi_{\rho}(g) $ is $\chi(g):\mathrm{Irr}(G) \to \C, \, \mu \mapsto 
\chi_{\mu}(g)$ (a column of the character table).
\end{prop}

Since $\chi_{\rho}(g), g\in G$ is the sum of eigenvalues of $\rho(g)$, all of 
which are roots of unity, we have:
\begin{coro}\label{cor:spectral_radius_adj_matrix}
 The eigenvalue of $A_{\Gamma(G, \rho)}$ with the maximal absolute value (the 
{\it spectral radius} of $A_{\Gamma(G, \rho)}$) is $\chi_\rho(1_G) = \dim \rho$. 
A corresponding eigenvector is then $$\dim:\mathrm{Irr}(G) \to \C, \, \mu \mapsto 
\dim(\mu).$$
\end{coro}

The spectral radius of a real matrix $A$ with non-negative entries has an important property given by the Frobenius-Perron theorem: the corresponding eigenspace of $A$ contains a vector all of whose entries are non-negative. 

In fact, if $A$ is an ``irreducible'' matrix (one which cannot be presented as a block matrix with non-trivial blocks), the spectral radius has multiplicity $1$ as an eigenvalue of $A$, and it is the only eigenvalue of $A$ for which an eigenvector with positive entries exists. For example, the adjacency matrix of a {\it strongly connected} graph $\Gamma$ (one where there is a directed walk from each vertex to any other vertex) is always irreducible, and hence its spectral radius has multiplicity $1$.


\subsection{Shapes of McKay graphs} Let us give some basic examples of graphs which can or cannot appear as McKay 
graphs $\Gamma(G, \rho)$ for some $G, \rho$.
\begin{exam}
\begin{enumerate}
\item Given a group $G$ and a representation $\rho$ of $G$, the representation 
$\rho$ is $1$-dimensional if and only if the McKay graph $\Gamma(G, \rho)$ is a 
disjoint union of directed cycles (some of the cycles might consist 
of $1$ vertex only, with a single loop). 

Any disjoint union of directed cycles of equal size can appear as a McKay 
graph: 
if we have $k$ disjoint directed cycles of size $n$, this is the 
McKay graph for $(\Z/n\Z \times G, \tau\boxtimes \triv)$ where $\tau$ is an 
irreducible representation of $\Z/n\Z$, and $G$ is any finite group with 
precisely $k$ irreducible representations (e.g. $G=\Z/k\Z$).
 \item  A McKay graph $\Gamma(G, \rho)$ cannot be a directed tree or a 
disjoint union of such: indeed, one cannot have ``leaves'' (i.e. vertices with 
only one edge coming in or going out) in a directed McKay graph.
 \item Considering the opposite extreme situation, a complete (undirected, 
loop-less) graph on $n$ vertices - such a graph can be obtained as 
$\Gamma(\Z/n\Z, \rho)$. Here $\rho$ is the representation of $\Z/n\Z$ on 
$\C^n/\Span\{(1,\ldots, 1)\}$, where $\Z/n\Z$ acts on $\C^n$ by shifting cyclically the 
coordinates. 
 \item A complete graph on $n$ vertices ($n>2$) cannot be obtained as a McKay 
graph $\Gamma(G, \rho)$ if we require $\rho \in \mathrm{Irr}(G)$. This is due to the 
fact that in any McKay graph $\Gamma(G, \rho)$ with $\rho \in \mathrm{Irr}(G)$, the 
vertex $\triv$ has only one neighbor: the vertex~$\rho$.
 
 Yet a given a complete (undirected, loop-less) graph on $n$ vertices, we can 
always find a McKay graph $\Gamma(G, \rho)$  containing our complete graph as a 
sub-graph: for instance, take $G = S_N$ (the symmetric group) for $N>>n$ and 
$\rho$ to be the reflection representation.
\end{enumerate}

\end{exam}


\subsection{Connected components in McKay graphs}

Given two vertices $\mu, \lam$ in a McKay graph, we can ask whether there is a directed walk from $\mu$ to $\lam$, or whether there is a walk from $\mu$ to $\lam$ if we forget about the directions of the edges in our graph. Although for general graphs these relations do not coincide, it turns out that in McKay graphs they do (see \cite[Section 3]{Browne}), so the notion of a connected component in a McKay graph is intuitive; each such component is ``strongly connected''.

Given a McKay graph $\Gamma(G, \rho)$, the connected component $\Gamma_{\triv}$ containing the vertex $\triv$ is called the {\it principal component}.

The following statement is well-known, see for example \cite[Proposition 1]{McKay}, \cite[Proposition 3.3]{Browne}:
\begin{lemma}\label{lem:connectivity}
 Let $\rho$ be a representation of the group $G$. Then $\rho$ is faithful if and 
only if $\Gamma(G, \rho)$ is connected (that is, for any $\mu, \lam \in \mathrm{Irr}(G)$ 
there exists a walk from $\mu$ to $\lam$ in $\Gamma(G, \rho)$).
\end{lemma}


In fact, given any representation $\rho$ of $G$ with kernel $N = \Ker(\rho)$, the principal component $\Gamma_{\triv}$ is isomorphic to $\Gamma(G/N, \rho)$. 
The connected components of $\Gamma(G, \rho)$ correspond to blocks in the adjacency matrix $A_{\Gamma(G, \rho)}$; the number of such blocks is then the multiplicity of the largest eigenvalue, implying the following lemma (see also \cite[Proposition 3.10, Lemma 4.1]{Browne}):
\begin{lemma}\label{lem:connected_comp_Browne} 
\mbox{}
 \begin{enumerate}
  \item The number of connected components of $\Gamma(G, \rho)$ is precisely the number of $G$-conjugacy classes of $N$; the latter equals the number of $G$-orbits in $\mathrm{Irr}({N})$ ($G$ acts on $\mathrm{Irr}(N)$ twisting the action of $N$ on the representations).
  
  In fact, there is a bijection
  $$\{\text{ connected components of } \Gamma(G, \rho)\} \longrightarrow \mathrm{Irr}(N)//G$$
  given as follows: to every connected component $\Gamma'$ of $\Gamma(G, \rho)$ corresponds a unique $G$-orbit $T \subset \mathrm{Irr}(N)$ such that $[\mu \downarrow^G_N: \tau] \neq 0$ for some $\mu \in X(\Gamma_T) \subset \mathrm{Irr}(G)$ and some $\tau \in T \subset N$.
  \item The group $\widehat{G}$ of characters $G \to \C^{\times}$ acts on the graph $\Gamma(G, \rho)$ by tensoring each vertex with the appropriate $1$-dimensional representation.
 \end{enumerate}
 
 
\end{lemma}

\begin{coro}\label{cor:conn_comp_isomorphic}
 Let $\Gamma'$ be a connected component of $\Gamma(G, \rho)$ with a vertex $\mu$ such that $\dim \mu=1$. Then the graph $\Gamma'$ is isomorphic to the principal component $\Gamma_{\triv}$.
\end{coro}


In fact, we have the following useful identity:
\begin{lemma}\label{lem:conn_comp_sum_of_sq}
  Let $T \in \mathrm{Irr}(N)//G$ and $\Gamma_T$ the corresponding connected component in $\Gamma(G, \rho)$. Let $\tau \in T$. We have:
 $$\sum_{\mu \in X(\Gamma_T)}(\dim \mu)^2 = \abs{G/N}\cdot (\dim\tau)^2 \cdot \abs{T}. $$
\end{lemma}
\begin{proof}

Given $\mu\in \mathrm{Irr}(G)$, the multiplicity $[\mu \downarrow^G_N: \tau]$ does not depend on choice of $\tau\in T$ by Clifford's theorem, and $[\mu \downarrow^G_N: \tau']=0$ for all $\tau'\notin T$. Hence, we have: $ \dim \mu = \dim \tau \cdot \abs{T} \cdot [\mu \downarrow^G_N: \tau] $ for $\mu \in X(\Gamma_T)$. In other words, $${[\tau\uparrow^G_N: \mu]}=[\mu \downarrow^G_N: \tau] = \frac{\dim \mu }{\dim \tau \cdot \abs{T}}.$$ 

Consider $\tau\uparrow^G_N$; by Lemma \ref{lem:connected_comp_Browne}, all its irreducible direct summands are in $X(\Gamma_T)$. We have: 
\begin{align*}
 \abs{G/N}\cdot\dim \tau  = \dim (\tau\uparrow^G_N) = \sum_{\mu \in X(\Gamma_T)}[\tau\uparrow^G_N: \mu]\dim \mu  = \sum_{\mu \in X(\Gamma_T)}\frac{(\dim \mu )^2}{\dim \tau  \cdot \abs{T}}
\end{align*}
The required statement is proved.
 
\end{proof}

\section{Auxiliary results}\label{sec:Aux}



\subsection{McKay classification}\label{ssec:mckay_class}
In this paper, the term {\it binary polyhedral group} (denoted by appending $\mathrm{B}$ to the notation of the 
group) stands for a group $G$ which is a double cover
of the polyhedral group of given type. 

{The polyhedral groups we consider are finite subgroups of $SO(3, \mathbb{R})$: these the rotational symmetry groups of the regular polyhedra, denoted by $\mathbf{T}$ (tetrahedral group), $\mathbf{O}$ (octahedral group) and $\mathbf{I}$ (icosahedral group), as well as the dihedral groups $\mathrm{Dih}_n$, $n\geq 2$. The latter are embedded naturally into $SO(3, \mathbb{R})$ by considering them as symmetries of ``degenerate" polyhedra: polygons which lie in the plane $XY$ of $\mathbb{R}^3$.

\begin{exam}\label{ex:Dih_2}
  The subgroup of $SO(3, \mathbb{R})$ corresponding to $\mathrm{Dih}_2 \cong C_2 \times C_2$ consists of matrices $ \begin{bmatrix}
  \pm 1 &0\\
  0 &\pm 1
  \end{bmatrix}$. These are precisely the rotations around the coordinate axes by multiples of $\pi$ radians.
\end{exam}}

\begin{exam}
The 
{\it binary dihedral group} (also known as ``dicyclic group'') $\mathrm{B}\mathrm{Dih}_{n}$ ($n\geq 2$) has presentation
$$\langle a, x | a^{2n} = 1, x^2=a^n, x^{-1}ax = a^{-1}\rangle.$$ It has center
$Z =\{1, a^n\}\cong C_2$ and the quotient $\mathrm{B}\mathrm{Dih}_{n}/Z$ is the dihedral group 
$\mathrm{Dih}_{n}$. For $n=2$, we have $\mathrm{Dih}_{2} \cong C_2 \times C_2$ and we obtain an 
isomorphism $\mathrm{B}\mathrm{Dih}_{2} \cong Q_8$ (the quaternion group).
\end{exam}

The following theorem is the most celebrated use of McKay graphs (see \cite[Proposition 4]{McKay} and~\cite{Steinberg}):
\begin{theorem}[McKay's theorem]\label{thrm:mckay_class}
The following list describes all McKay graphs for pairs $(G, \rho)$ 
where $G\subset SU(2)$ is a finite group and $\rho$ is the respective
$2$-dimensional representation of $G$.

\end{theorem}


\begin{center}
\begin{tabular}{  !{\color{blue}\vrule width 2pt} c|c !{\color{blue}\vrule 
width 2pt} } 
 \specialrule{.2em}{.0em}{.0em}
 \xrowht{20pt}
 Affine Dynkin diagram & Group $G$ \\ 
 \specialrule{.2em}{.0em}{.0em} &
 \\$\tilde A_n$, $n\geq 0$  &cyclic\\ $$
\begin{tikzcd}[arrows=-]
&&\bigstar\arrow[drr]&&\\
  \circled{1}\arrow[rru]\arrow[r] & \circled{1}\arrow[r] 
  &\cdots  \arrow[r] &
  \circled{1} \arrow[r] & \circled{1}
\end{tikzcd}
$$
   &$\Z/(n+1)\Z$ \\ &  size: $n+1$\\
 \hline  &\\
 $\tilde D_n$ , $n\geq 4$ & binary dihedral  \\ $$
\begin{tikzcd}[arrows=-]
\bigstar\arrow[rd]&&&& \circled{1}\\ 
   & \circled{2}\arrow[r] 
  &\cdots  \arrow[r] &
  \circled{2}  \arrow[ru]\arrow[rd] & \\
\circled{1}\arrow[ru]&&&&\circled{1}
\end{tikzcd}
$$
& $\mathrm{B}\mathrm{Dih}_{n-2}$ \\ & size: $4(n-2)$ \\ 
 \hline &\\
 $\tilde E_6$ & binary tetrahedral \\
 $$
\begin{tikzcd}[arrows=-]
    &&\bigstar\arrow[d]&&\\
    && \circled{2}\arrow[d]&&\\ 
  \circled{1}\arrow[r] & \circled{2}\arrow[r] 
  &\circled{3}  \arrow[r] & 
  \circled{2} \arrow[r]  & \circled{1} 
\end{tikzcd}
$$

  &$\mathrm{B}{\bf T}$ \\ & size: $24$ \\ 
 \hline & \\
 $\tilde E_7$ & binary octahedral\\
 $$
\begin{tikzcd}[sep=small,arrows=-]
    &&& \circled{2}\arrow[d]&&&\\ 
 \bigstar\arrow[r] & \circled{2}\arrow[r] & \circled{3}\arrow[r] 
  &\circled{4} \arrow[r] &
  \circled{3} \arrow[r]  & \circled{2}\arrow[r] & \circled{1}
\end{tikzcd}
$$

& $\mathrm{B}{\bf O}$ \\ & size: $48$   \\ 
 \hline  &\\
 $\tilde E_8$ & binary icosahedral\\
$$
\begin{tikzcd}[sep=small,arrows=-]
    &&&&&
  \circled{3}\arrow[d]&&\\ 
  \bigstar\arrow[r]&\circled{2}\arrow[r]&\circled{3}\arrow[r]&\circled{4}\arrow[r] 
& \circled{5}\arrow[r] & \circled{6}\arrow[r] 
  & \circled{4} \arrow[r] & \circled{2} 
\end{tikzcd}
$$
& $\mathrm{B}{\bf I}$ \\ & size: $120$ \\ 
 \specialrule{.2em}{.0em}{.0em}
\end{tabular}
\end{center}

The proof of Theorem \ref{thrm:mckay_class} relies on the classification of all finite undirected connected simply-laced graphs $\Gamma$ with spectral radius $2$. The classification of such graphs states that they must be affine Dynkin diagrams of types $A$, $D$ or $E$, {which was proved by Smith in \cite{Smith} (cf. \cite[Proposition~4]{McKay})}. 



\begin{prop}\label{prop:McKay_classification}
  Let $G$ be a finite group and $\rho$ be its self-dual faithful 
representation (so $\Gamma(G, \rho)$ is connected and undirected). Assume 
$\dim \rho = 2$. 

Then either 
$G\subset SU(2)$ (hence $G$ appears in the list of Theorem \ref{thrm:mckay_class} and $\Gamma(G, \rho)$ is of the form above), 
or $G=\mathrm{Dih}_n$ and $\rho$ is the 
complexification of the tautological $2$-dimensional representation of 
$\mathrm{Dih}_n$. In the latter case, $\Gamma(G, \rho)$ is drawn below.
\end{prop}


\begin{center}
\begin{tabular}{  !{\color{blue}\vrule width 2pt} c|c|c !{\color{blue}\vrule 
width 2pt} } 
 \specialrule{.2em}{.0em}{.0em}
 \xrowht{20pt}
 parity of $n$ & {\small McKay graph for $\mathrm{Dih}_n$, tautological $\rho$, $\dim(\rho)=2$} & number of vertices\\ 
 \specialrule{.2em}{.0em}{.0em}
 $n$ even  & $$
\begin{tikzcd}[arrows=-]
\bigstar\arrow[rd]&&&&\circled{1}\\ 
   & \circled{2}\arrow[r] 
  &\cdots  \arrow[r] &
  \circled{2}\arrow[ru]\arrow[rd] & \\ 
\circled{1}\arrow[ru]&&&&\circled{1} 
\end{tikzcd}
$$ &$\frac{n}{2}+3$
 \\
 \hline
 $n$ odd & $$
\begin{tikzcd}[arrows=-]
\bigstar\arrow[rd]&&&&{}\\
   & \circled{2}\arrow[r]
  &\cdots  \arrow[r] &
  \circled{2}  \arrow[loop right]{r} & \\ 
\circled{1}\arrow[ru]&&&&{} 
\end{tikzcd}
$$ &$\frac{n+3}{2}$\\ 
 \specialrule{.2em}{.0em}{.0em}
\end{tabular}
\end{center}

\begin{proof}
Since $\rho$ is 
self dual, it either has a symplectic $G$-invariant form (i.e. $\rho$ is of 
quaternionic type) or a symmetric $G$-invariant form (i.e. $\rho$ is of 
real type). In the former case, $G 
\subset SL(2)\cap U(2) =SU(2)$ and we are in the situation of Theorem~\ref{thrm:mckay_class}. In the latter case, the representation $(\rho, V)$ is 
of real type, the corresponding $2$-dimensional real representation 
$(\rho_{\mathbb R}, V_{\mathbb R})$ satisfies: $\rho \cong 
\C\otimes_{\mathbb R}\rho_{\mathbb R}$, where $\rho_{\mathbb R}$ is a faithful $2$-dimensional 
representation of $G$ over $\mathbb R$. Hence $G\subset {\rm Dih}_n$ for some $n\geq 1$. But in that case $G$ it either cyclic (and then we are again in the situation of Theorem 
\ref{thrm:mckay_class}) or $G$ is dihedral. This 
completes the proof of the proposition.
\end{proof}




